% GENERATED FILE. Edit canonical lessons, then run npm run generate:books.
% canonical-source-hash: fnv1a64:1a0dd710c094fd36
% canonical-lessons: GU-C431-godam, GU-C431-puravtho, GU-C431-unap, GU-C431-jaththabandh, GU-C431-sangharvu

\chapter{Stock and Shortages}
\label{ch:gu-stock-and-shortages}
\hlchaptermodality{431}

\begin{chapteropening}
\textbf{By the end of this chapter:} \emph{I can talk about supplies running short.}

Complete the last lesson of chapter 431: I can talk about supplies running short.
\end{chapteropening}

\section[\texorpdfstring{\gu{ગોદામ} (godām)}{ગોદામ (godām)}]{\gu{ગોદામ} (godām) — a warehouse; a godown}

\label{lesson:GU-C431-godam}



\begin{quote}
Before the new one: say the Gujarati for a regular buyer, then the Gujarati for to weigh out.
\end{quote}

\subsection*{You'll want to know: \gu{ગોદામ}}
\textbf{\gu{ગોદામ}} — \emph{godām} — \textquotedblleft{}a warehouse; a godown\textquotedblright{}.

\begin{quote}
\gu{વેપારીના} \gu{ગોદામમાં} \gu{ઘઉં} \gu{ભરેલા} \gu{છે}.
\end{quote}

\begin{quote}
\emph{vepārīnā godāmmā̃ ghaũ bharelā che.}
\end{quote}

\begin{quote}
The merchant's warehouse is full of wheat.
\end{quote}

A neuter noun.

\begin{cousinweb}
\textbf{\gu{ગોદામ}} (\emph{godām}) is usually traced to Malay \emph{gudang}, \textquotedblleft{}a storehouse\textquotedblright{}, which came round the coasts through Portuguese; English \emph{godown} is the same word.
\end{cousinweb}

\begin{grammarlens}[title={What you've built}]
One more word for getting around.
\end{grammarlens}

\subsection*{Your turn}
\begin{itemize}
\raggedright
  \item \emph{Say it:} \emph{godām}
  \item \emph{Say it:} \emph{godām}, once more
  \item \emph{Say it:} \emph{vepārīnā godāmmā̃ ghaũ bharelā che.} — The merchant's warehouse is full of wheat.
  \item \emph{Recall:} say the Gujarati for a regular buyer, then the Gujarati for to weigh out, then say \emph{godām} again
\end{itemize}

\begin{tcolorbox}[breakable,colback=teal!4,colframe=teal!35!black,title={Before you move on}]
What does \gu{ગોદામ} mean? (\textquotedblleft{}A warehouse; a godown\textquotedblright{}.) Say it once more.
\end{tcolorbox}
\section[\texorpdfstring{\gu{પુરવઠો} (puravṭho)}{પુરવઠો (puravṭho)}]{\gu{પુરવઠો} (puravṭho) — a supply (of goods, water or power)}

\label{lesson:GU-C431-puravtho}



\begin{quote}
Before the new one: say the Gujarati for to weigh out, then the Gujarati for a warehouse.
\end{quote}

\subsection*{You'll want to know: \gu{પુરવઠો}}
\textbf{\gu{પુરવઠો}} — \emph{puravṭho} — \textquotedblleft{}a supply (of goods, water or power)\textquotedblright{}.

\begin{quote}
\gu{ગામમાં} \gu{પાણીનો} \gu{પુરવઠો} \gu{ઓછો} \gu{છે}.
\end{quote}

\begin{quote}
\emph{gāmmā̃ pāṇīno puravṭho ocho che.}
\end{quote}

\begin{quote}
The water supply in the village is low.
\end{quote}

A masculine noun.

\begin{cousinweb}
\textbf{\gu{પુરવઠો}} (\emph{puravṭho}) seems to belong with the verb CDIAL lists for Gujarati, \emph{pūrvũ}, \textquotedblleft{}to fill\textquotedblright{}, and Marathi \emph{puraviṇẽ}, \textquotedblleft{}to supply to the full\textquotedblright{}, under Sanskrit \emph{pūrayati}, \textquotedblleft{}fills\textquotedblright{} (CDIAL 8335); CDIAL gives no form for \gu{પુરવઠો} itself, so the link is hedged.
\end{cousinweb}

\begin{grammarlens}[title={What you've built}]
One more word for getting around.
\end{grammarlens}

\subsection*{Your turn}
\begin{itemize}
\raggedright
  \item \emph{Say it:} \emph{puravṭho}
  \item \emph{Say it:} \emph{puravṭho}, once more
  \item \emph{Say it:} \emph{gāmmā̃ pāṇīno puravṭho ocho che.} — The water supply in the village is low.
  \item \emph{Recall:} say the Gujarati for to weigh out, then the Gujarati for a warehouse, then say \emph{puravṭho} again
\end{itemize}

\begin{tcolorbox}[breakable,colback=teal!4,colframe=teal!35!black,title={Before you move on}]
What does \gu{પુરવઠો} mean? (\textquotedblleft{}A supply (of goods, water or power)\textquotedblright{}.) Say it once more.
\end{tcolorbox}
\section[\texorpdfstring{\gu{ઊણપ} (ūṇap)}{ઊણપ (ūṇap)}]{\gu{ઊણપ} (ūṇap) — a shortfall; a lack}

\label{lesson:GU-C431-unap}



\begin{quote}
Before the new one: say the Gujarati for a warehouse, then the Gujarati for a supply.
\end{quote}

\subsection*{You'll want to know: \gu{ઊણપ}}
\textbf{\gu{ઊણપ}} — \emph{ūṇap} — \textquotedblleft{}a shortfall; a lack\textquotedblright{}.

\begin{quote}
\gu{ઉનાળામાં} \gu{ગામમાં} \gu{પાણીની} \gu{ઊણપ} \gu{રહે} \gu{છે}.
\end{quote}

\begin{quote}
\emph{unāḷāmā̃ gāmmā̃ pāṇīnī ūṇap rahe che.}
\end{quote}

\begin{quote}
In summer the village runs short of water.
\end{quote}

A feminine noun.

\begin{cousinweb}
\textbf{\gu{ઊણપ}} (\emph{ūṇap}) is built on \gu{ઊણું}, \textquotedblleft{}wanting, short\textquotedblright{}, which CDIAL lists under Sanskrit \emph{ūna}, \textquotedblleft{}deficient\textquotedblright{} (CDIAL 2405): what is missing from the full amount.
\end{cousinweb}

\begin{grammarlens}[title={What you've built}]
One more word for getting around.
\end{grammarlens}

\subsection*{Your turn}
\begin{itemize}
\raggedright
  \item \emph{Say it:} \emph{ūṇap}
  \item \emph{Say it:} \emph{ūṇap}, once more
  \item \emph{Say it:} \emph{unāḷāmā̃ gāmmā̃ pāṇīnī ūṇap rahe che.} — In summer the village runs short of water.
  \item \emph{Recall:} say the Gujarati for a warehouse, then the Gujarati for a supply, then say \emph{ūṇap} again
\end{itemize}

\begin{tcolorbox}[breakable,colback=teal!4,colframe=teal!35!black,title={Before you move on}]
What does \gu{ઊણપ} mean? (\textquotedblleft{}A shortfall; a lack\textquotedblright{}.) Say it once more.
\end{tcolorbox}
\section[\texorpdfstring{\gu{જથ્થાબંધ} (jaththābandh)}{જથ્થાબંધ (jaththābandh)}]{\gu{જથ્થાબંધ} (jaththābandh) — wholesale; in bulk}

\label{lesson:GU-C431-jaththabandh}



\begin{quote}
Before the new one: say the Gujarati for a supply, then the Gujarati for a shortfall.
\end{quote}

\subsection*{You'll want to know: \gu{જથ્થાબંધ}}
\textbf{\gu{જથ્થાબંધ}} — \emph{jaththābandh} — \textquotedblleft{}wholesale; in bulk\textquotedblright{}.

\begin{quote}
\gu{વેપારી} \gu{જથ્થાબંધ} \gu{ભાવે} \gu{માલ} \gu{વેચે} \gu{છે}.
\end{quote}

\begin{quote}
\emph{vepārī jaththābandh bhāve māl vece che.}
\end{quote}

\begin{quote}
The trader sells goods at wholesale prices.
\end{quote}

\begin{cousinweb}
\textbf{\gu{જથ્થાબંધ}} (\emph{jaththābandh}) is built from \gu{જથ્થો}, \textquotedblleft{}a quantity, a bulk\textquotedblright{}, and -\gu{બંધ}, \textquotedblleft{}bound, tied\textquotedblright{}, usually traced to Persian \emph{band}: goods sold tied up in bulk.
\end{cousinweb}

\begin{grammarlens}[title={What you've built}]
One more word for getting around.
\end{grammarlens}

\subsection*{Your turn}
\begin{itemize}
\raggedright
  \item \emph{Say it:} \emph{jaththābandh}
  \item \emph{Say it:} \emph{jaththābandh}, once more
  \item \emph{Say it:} \emph{vepārī jaththābandh bhāve māl vece che.} — The trader sells goods at wholesale prices.
  \item \emph{Recall:} say the Gujarati for a supply, then the Gujarati for a shortfall, then say \emph{jaththābandh} again
\end{itemize}

\begin{tcolorbox}[breakable,colback=teal!4,colframe=teal!35!black,title={Before you move on}]
What does \gu{જથ્થાબંધ} mean? (\textquotedblleft{}Wholesale; in bulk\textquotedblright{}.) Say it once more.
\end{tcolorbox}
\section[\texorpdfstring{\gu{સંઘરવું} (sangharvũ)}{સંઘરવું (sangharvũ)}]{\gu{સંઘરવું} (sangharvũ) — to store up; to hoard}

\label{lesson:GU-C431-sangharvu}



\begin{quote}
Before the new one: say the Gujarati for a shortfall, then the Gujarati for wholesale.
\end{quote}

\subsection*{You'll want to know: \gu{સંઘરવું}}
\textbf{\gu{સંઘરવું}} — \emph{sangharvũ} — \textquotedblleft{}to store up; to hoard\textquotedblright{}.

\begin{quote}
\gu{વેપારીએ} \gu{ગોદામમાં} \gu{ખાંડ} \gu{સંઘરી}.
\end{quote}

\begin{quote}
\emph{vepārīe godāmmā̃ khā̃ḍ sanghrī.}
\end{quote}

\begin{quote}
The trader stored up sugar in his warehouse.
\end{quote}

Past, with \gu{તેણે} (he or she …): \gu{તેણે} \gu{ઘઉં} \gu{સંઘર્યા}, \emph{teṇe ghaũ sangharyā}.

\begin{cousinweb}
\textbf{\gu{સંઘરવું}} (\emph{sangharvũ}) is built on Sanskrit \emph{saṁgraha}, \textquotedblleft{}holding together, gathering in\textquotedblright{} (Monier-Williams): to keep gathering things and not let them go.
\end{cousinweb}

\begin{grammarlens}[title={What you've built}]
One more word for getting around.
\end{grammarlens}

\subsection*{Your turn}
\begin{itemize}
\raggedright
  \item \emph{Say it:} \emph{sangharvũ}
  \item \emph{Say it:} \emph{sangharvũ}, once more
  \item \emph{Say it:} \emph{vepārīe godāmmā̃ khā̃ḍ sanghrī.} — The trader stored up sugar in his warehouse.
  \item \emph{Recall:} say the Gujarati for a shortfall, then the Gujarati for wholesale, then say \emph{sangharvũ} again
\end{itemize}

\begin{tcolorbox}[breakable,colback=teal!4,colframe=teal!35!black,title={Before you move on}]
What does \gu{સંઘરવું} mean? (\textquotedblleft{}To store up; to hoard\textquotedblright{}.) Say it once more.
\end{tcolorbox}
